\section{Time derivatives of the actual signed third growth}
\label{tgdc:section}

This calculation is on the actual interval already constructed in
\emph{The actual third growth target on the moving second hold} and
\emph{Exact third growth on the moving second holding state}. It gives
quantitative derivatives of that trajectory for the source-frame map.
All derivatives in this section are with respect to coupled physical
time $t_c$, and $\tau=t_c-t_b$ denotes elapsed coupled time. The notation
does not use the cylindrical spatial radius for elapsed time.

\subsection{Original data, evolving parent, and proved interval}

Retain the source parameters
\begin{equation}\label{tgdc:original}
\begin{gathered}
 \nu=\sqrt{A_0},\qquad
 \mu=\lambda_0=\lceil\sqrt{A_0}/2\rceil,\qquad
 K_j=\mu\widehat\lambda_j,\\
 S_3=\frac{A_0}{\mu}\widehat\lambda_3^{-k_3-6},\qquad
 P_3^*=\frac{A_0}{\mu}\widehat\lambda_3^{-7/8}=S_3e^{L_3}.
\end{gathered}
\end{equation}
The physical common diffusivity is $d>0$; it is distinct from $\nu$.
Use the actual second return $t_b$, its values $h_b,A_{1,b},P_{2,b}$,
and the original angles $s_*,s_2,s_3$. Put
\begin{equation}\label{tgdc:constants}
\begin{gathered}
 a=\sin s_2,\quad \chi_b=\sqrt{h_b^2+a^2},\quad
 \beta_b=\arctan(a/h_b),\quad
 b=\sin(s_3+\beta_b),\quad g_b=\cos(s_3+\beta_b),\\
 c=bh_b-ag_b=\chi_b\sin s_3,\qquad
 S_*=A_0\sin s_*,\quad C_*=A_0\cos s_*.
\end{gathered}
\end{equation}
Here $a,b,c,S_*,C_*>0$, and the letter $b$ is the original oriented
geometric constant, whereas $b_3$ below is an amplitude coefficient.
The moving parent is exactly
\begin{equation}\label{tgdc:parent}
\begin{aligned}
 \chi^2&=h^2+a^2,& \dot h&=a\Omega_1,&
 \dot\Omega_1&=aA_1/\chi-dK_1^2\Omega_1,\\
 \dot A_1&=-S_*\Omega_1-dK_1^2A_1,&
 \dot P_2&=-dK_2^2\chi^2P_2,&
 \dot g&=b\Omega_1,\\
 g&=g_b+(b/a)(h-h_b),&
 N_3&=b(A_1+C_*)+gS_*+cK_2P_2,&
 R_3&=g^2+b^2.
\end{aligned}
\end{equation}
The retained feedback is $\dot\alpha=-a^2\Omega_1/\chi^2$ and
$\Omega_2=0$. In particular the older temperature and first vorticity
are not frozen. The cancellation in the complete numerator is
\begin{equation}\label{tgdc:numerator}
 \dot N_3=-dE_3,\qquad
 E_3=bK_1^2A_1+cK_2^3\chi^2P_2.
\end{equation}

Write $N_i=N_3(t_b)$, $a_i=N_i/K_3$,
$b_i=K_3\sin s_3$, $\gamma_i=\sqrt{a_i b_i}$ and
$u_i=\sqrt{b_i/a_i}$. The exact new equations, including their
original signed seed, are
\begin{equation}\label{tgdc:actual_pair}
\begin{gathered}
 y=\binom{\Theta_3}{\Omega_3},\qquad
 \dot y=C_3y,\qquad
 C_3=\begin{pmatrix}-\lambda_3&a_3\\b_3&-\lambda_3\end{pmatrix},\\
 a_3=\frac{N_3}{K_3R_3},\qquad
 b_3=\frac{K_3c}{\chi},\qquad
 \lambda_3=dK_3^2R_3,\qquad
 y(t_b)=-S_3\binom1{u_i}.
\end{gathered}
\end{equation}
There are no new amplitude controls in this growth pair:
$c_\vartheta=c_\omega=0$. The activation of the compact wave is a
separate original factor $\eta(\Gamma_3\tau)$; it does not replace
\eqref{tgdc:actual_pair} by a forced coefficient trajectory.

For clarity, the following facts about this actual interval were proved
in the third-threshold construction. They identify the trajectory to
which the estimates below apply:
\begin{equation}\label{tgdc:proved_boxes}
\begin{gathered}
 I_3=[t_b,t_3^a],\quad t_3^a-t_b=\tau_a<14/\sigma_1,\quad
 P_3=-\Theta_3>0,\quad P_3(t_b)=S_3,\quad P_3(t_3^a)=P_3^*,\\
 u=\Omega_3/\Theta_3,\qquad
 \frac{u_i}{\sqrt2}\leq u\leq2u_i,\qquad
 \frac{a_i}4\leq a_3\leq a_i,\qquad
 \frac{b_i}2\leq b_3\leq b_i,\\
 1\leq R_3\leq\overline R:=\left(\frac{265}{256}\right)^2<2,
 \quad 1\leq\frac\chi{\chi_b}\leq\frac{259}{256},\quad
 0<dK_3^2\leq\frac{\gamma_i}{16\sqrt2},\\
 \frac{\gamma_i}{8\sqrt2}\leq (\log P_3)'\leq2\gamma_i,
 \quad 0<A_1\leq A_{1,b},\quad 0<P_2\leq P_{2,b},\quad
 0<\Omega_1<9\sigma_1,\\
 h_b\leq h\leq h_b+144a,\qquad
 g_b\leq g\leq g_b+144b,\qquad
 \tfrac12N_i\leq N_3\leq N_i,\qquad 1\leq\chi^2<10.
\end{gathered}
\end{equation}
The value of $d$ is precisely the previously selected range
$0<d\leq d_{3,\max}$, retaining its $K_3^{-2}$ dependence.
Thus \eqref{tgdc:proved_boxes} does not assert a new unrestricted
viscosity range or assume a further return.

\subsection{Signed amplitudes and first derivatives with the diffusion retained}

\begin{theorem}\label{tgdc:first_derivatives}
For every $t_c\in I_3$, both $\Theta_3(t_c)$ and $\Omega_3(t_c)$
are negative and strictly decreasing. Their exact component rates are
\begin{equation}\label{tgdc:exact_first}
 \dot\Theta_3=-P_3(a_3u-\lambda_3),\qquad
 \dot\Omega_3=-P_3(b_3-\lambda_3u).
\end{equation}
They obey the bounds
\begin{equation}\label{tgdc:component_bounds}
\begin{gathered}
 S_3e^{\gamma_i\tau/(8\sqrt2)}
 \leq P_3(t_c)\leq\min\{P_3^*,S_3e^{2\gamma_i\tau}\},\\
 \frac{u_i}{\sqrt2}P_3\leq |\Omega_3|\leq2u_iP_3,\\
 \left(\frac{\gamma_i}{4\sqrt2}-dK_3^2\overline R\right)P_3
 \leq |\dot\Theta_3|\leq2\gamma_iP_3,\\
 u_i\left(\frac{\gamma_i}{2}-2dK_3^2\overline R\right)P_3
 \leq |\dot\Omega_3|\leq\gamma_i u_iP_3.
\end{gathered}
\end{equation}
In particular both lower derivative bounds are positive, and
\begin{equation}\label{tgdc:norm_bounds}
\begin{gathered}
 P_3\sqrt{1+u_i^2/2}\leq |y|_2
       \leq P_3\sqrt{1+4u_i^2}\leq P_3^*\sqrt{1+4u_i^2},\\
 |\dot y|_2\leq\gamma_iP_3\sqrt{4+u_i^2}
       \leq\gamma_iP_3^*\sqrt{4+u_i^2}.
\end{gathered}
\end{equation}
The original beginning and attained endpoint have the exact derivatives
\begin{equation}\label{tgdc:endpoint_derivatives}
\begin{aligned}
 \dot y(t_b)&=-S_3(\gamma_i-dK_3^2)\binom1{u_i},\\
 y(t_3^a)&=-P_3^*\binom1{u(t_3^a)},\\
 \dot y(t_3^a)&=-P_3^*
 \binom{a_3(t_3^a)u(t_3^a)-dK_3^2R_3(t_3^a)}
       {b_3(t_3^a)-dK_3^2R_3(t_3^a)u(t_3^a)}.
\end{aligned}
\end{equation}
\end{theorem}
\begin{proof}
The signs and quotient in \eqref{tgdc:proved_boxes} give
$y=-P_3(1,u)^T$. Substituting this expression in the exact pair proves
\eqref{tgdc:exact_first}. Integrating the already proved logarithmic
rate from $t_b$ and using the attained endpoint and strict increase
of $P_3$ gives the first line of \eqref{tgdc:component_bounds}.
The quotient box gives its second line. Since $a_i u_i=\gamma_i$
and $b_i=\gamma_i u_i$, direct multiplication of the positive boxes
gives
\[
 \frac{\gamma_i}{4\sqrt2}\leq a_3u\leq2\gamma_i,
 \qquad
 \frac{\gamma_i u_i}{2}\leq b_3\leq\gamma_i u_i,
 \qquad 0<\lambda_3\leq dK_3^2\overline R.
\]
Consequently the displayed lower bound for $a_3u-\lambda_3$ follows.
It is at least $\gamma_i/(8\sqrt2)>0$, because
$\overline R<2$ and $dK_3^2\leq\gamma_i/(16\sqrt2)$.
Similarly,
\[
 b_3-\lambda_3u
 \geq u_i(\gamma_i/2-2dK_3^2\overline R)
 >\gamma_i u_i(1/2-1/(4\sqrt2))
 >\gamma_i u_i/4>0.
\]
Thus both derivatives in \eqref{tgdc:exact_first} are strictly negative.
Dropping their nonnegative diffusion losses only for the purpose of
an upper bound gives the two upper derivative estimates; the exact
rates continue to contain the losses. Squaring the two amplitude and
two derivative bounds, adding them, and taking the positive square
root proves \eqref{tgdc:norm_bounds}. Finally $R_3(t_b)=1$,
$u(t_b)=u_i$, $a_3(t_b)=a_i$ and $b_3(t_b)=b_i$ give the first
endpoint identity. Substitution of the proved target value gives
the other two. In particular no later endpoint was substituted for
the actual $u(t_3^a)$.
\end{proof}

\subsection{Explicit coefficient rates}

Set $\overline H=h_b+144a$, $\overline G=g_b+144b$ and
$\overline W=9\sigma_1$. Differentiating the complete coefficients
gives
\begin{equation}\label{tgdc:coefficient_rates}
\begin{aligned}
 \dot a_3&=-a_3\left(\frac{dE_3}{N_3}
                            +\frac{2bg\Omega_1}{R_3}\right),\\
 \dot b_3&=-b_3\frac{ah\Omega_1}{\chi^2},\qquad
 \dot\lambda_3=2dK_3^2bg\Omega_1.
\end{aligned}
\end{equation}
Every quantity on the right comes from the moving parent. In particular
the exact numerator derivative \eqref{tgdc:numerator} is used before
any upper estimate. Because $K_1<K_2$, $\chi^2<10$ and all summands
of $N_3$ are positive, $E_3\leq10K_2^2N_3$. Thus
\begin{equation}\label{tgdc:coefficient_rate_bounds}
\begin{gathered}
 |\dot a_3|\leq
 a_i(10dK_2^2+2b\overline G\overline W),\qquad
 |\dot b_3|\leq
 b_i a\overline H\overline W\chi_b^{-2},\\
 |\dot\lambda_3|\leq
 2dK_3^2 b\overline G\overline W.
\end{gathered}
\end{equation}
These formulas also display the distinct inherited $K_2^2$ loss and
new $K_3^2$ loss. They do not replace either by the source scaling
parameter $\nu$.

\subsection{A closed all-orders derivative calculation}

Here are finite recurrences whose values are explicit upper bounds
from the actual birth data and the already proved interval. They use
ordinary derivatives, with no change to physical time or amplitudes.
For integers $n\geq0$ and $m\geq2$ define the product expression
\[
 \mathcal P_n(X^{(1)},\ldots,X^{(m)})
 =\sum_{j_1+\cdots+j_m=n}
       \frac{n!}{j_1!\cdots j_m!}
       X^{(1)}_{j_1}\cdots X^{(m)}_{j_m}.
\]
This notation is a finite sum, and repeated arguments mean repeated
factors. Define the nonnegative arrays $H,W,A,P,Q,G,N,J$ initially by
\begin{equation}\label{tgdc:majorant_initial}
\begin{gathered}
 H_0=\overline H,\quad W_0=\overline W,\quad
 A_0^{\rm bd}=A_{1,b},\quad P_0=P_{2,b},\quad
 Q_0=\chi_b^{-1},\quad G_0=\overline G,\quad
 N_0=N_i,\quad J_0=1.
\end{gathered}
\end{equation}
In the following recurrences the array $A$ has entries $A_n^{\rm bd}$;
the superscript distinguishes its zeroth entry from the unchanged
physical constant $A_0$ in \eqref{tgdc:original}. Write
\[
 X_n=\mathcal P_n(H,H)+a^2\mathbf1_{n=0}.
\]
Having computed all eight arrays up to index $n$, compute index $n+1$
by the explicit formulas
\begin{equation}\label{tgdc:parent_majorants}
\begin{aligned}
 H_{n+1}&=aW_n,\\
 W_{n+1}&=a\mathcal P_n(A,Q)+dK_1^2W_n,\\
 A_{n+1}^{\rm bd}&=S_*W_n+dK_1^2A_n^{\rm bd},\\
 P_{n+1}&=dK_2^2\mathcal P_n(X,P),\\
 Q_{n+1}&=a\mathcal P_n(H,W,Q,Q,Q),\\
 G_{n+1}&=bW_n,\\
 N_{n+1}&=dbK_1^2A_n^{\rm bd}
                 +dcK_2^3\mathcal P_n(X,P),\\
 J_{n+1}&=2b\mathcal P_n(G,W,J,J).
\end{aligned}
\end{equation}
For the phase-length array use
\[
 R_0^{\rm bd}=\overline R,\qquad
 R_n^{\rm bd}=\mathcal P_n(G,G)\quad(n\geq1).
\]
The coefficient arrays are
\begin{equation}\label{tgdc:coefficient_majorants}
\begin{gathered}
 \mathfrak A_0=a_i,\quad \mathfrak B_0=b_i,\quad
 \mathfrak D_0=dK_3^2\overline R,\\
 \mathfrak A_n=K_3^{-1}\mathcal P_n(N,J),\qquad
 \mathfrak B_n=K_3cQ_n,\qquad
 \mathfrak D_n=dK_3^2R_n^{\rm bd}\quad(n\geq1).
\end{gathered}
\end{equation}
No optimization of these finite sums is needed for their validity.
In particular $\mathfrak A_1$ may be replaced by the smaller of its
value above and the first bound in \eqref{tgdc:coefficient_rate_bounds},
and similarly for the other two first coefficient derivatives.

The amplitude arrays start with the actual target and the sharper
first-derivative bounds already proved:
\begin{equation}\label{tgdc:amplitude_majorants}
\begin{gathered}
 T_0=P_3^*,\qquad O_0=2u_iP_3^*,\qquad
 T_1=2\gamma_iP_3^*,\qquad O_1=\gamma_i u_iP_3^*,\\
 T_{n+1}=\sum_{j=0}^n\binom nj
      (\mathfrak D_jT_{n-j}+\mathfrak A_jO_{n-j}),\qquad n\geq1,\\
 O_{n+1}=\sum_{j=0}^n\binom nj
      (\mathfrak B_jT_{n-j}+\mathfrak D_jO_{n-j}),\qquad n\geq1.
\end{gathered}
\end{equation}
The choice to begin these recurrences at $n=1$ preserves the stronger
first-derivative estimates instead of substituting a larger matrix-norm
bound.

\begin{theorem}\label{tgdc:all_orders}
For every integer $n\geq0$ the finite arrays above satisfy on $I_3$
\begin{equation}\label{tgdc:all_order_bounds}
\begin{gathered}
 |\partial_{t_c}^n\Theta_3|\leq T_n,\qquad
 |\partial_{t_c}^n\Omega_3|\leq O_n,\qquad
 |\partial_{t_c}^ny|_2\leq\sqrt{T_n^2+O_n^2},\\
 |\partial_{t_c}^na_3|\leq\mathfrak A_n,\qquad
 |\partial_{t_c}^nb_3|\leq\mathfrak B_n,\qquad
 |\partial_{t_c}^n\lambda_3|\leq\mathfrak D_n.
\end{gathered}
\end{equation}
The exact, signed amplitude derivative recurrence underlying these
bounds is
\begin{equation}\label{tgdc:exact_derivative_recurrence}
 \partial_{t_c}^{n+1}y
     =\sum_{j=0}^n\binom nj
            (\partial_{t_c}^jC_3)(\partial_{t_c}^{n-j}y),
 \qquad n\geq0.
\end{equation}
In particular the exact second derivative contains the changing
coefficient and the complete diffusion square:
\begin{equation}\label{tgdc:second_derivative}
 \ddot y=
 \begin{pmatrix}
 \lambda_3^2+a_3b_3-\dot\lambda_3&\dot a_3-2\lambda_3a_3\\
 \dot b_3-2\lambda_3b_3&\lambda_3^2+a_3b_3-\dot\lambda_3
 \end{pmatrix}y.
\end{equation}
\end{theorem}
\begin{proof}
Introduce the exact auxiliary functions $q=\chi^{-1}$ and
$j=R_3^{-1}$, both well-defined by \eqref{tgdc:proved_boxes}.
Direct differentiation of their specified inverses gives
\[
 \dot q=-ah\Omega_1q^3,\qquad
 \dot j=-2bg\Omega_1j^2.
\]
These identities retain the actual phase lengths; the symbols do not
replace them by constant values. The complete equations for
$(h,\Omega_1,A_1,P_2,q,g,N_3,j)$ are therefore the signed polynomial
system consisting of \eqref{tgdc:parent}, \eqref{tgdc:numerator} and
these two equations, with $\chi^2=h^2+a^2$ in the parent losses.
Their zero-order absolute values are bounded respectively by
$(H_0,W_0,A_0^{\rm bd},P_0,Q_0,G_0,N_0,J_0)$:
the claims for $q$ and $j$ use $\chi\geq\chi_b$ and $R_3\geq1$.

For any $m$ differentiable factors the ordinary $n$th derivative of
their product is $\mathcal P_n$ applied to their derivative arrays.
To prove this formula without assuming it, start with $n=0$.
Differentiating a term of order $n$ once chooses one factor whose
index increases. At indices summing to $n+1$, the resulting coefficient
is the sum of $n!/(j_1!\cdots(j_k-1)!\cdots j_m!)$ over indices
$k$ with $j_k>0$. It is
$n!(j_1+\cdots+j_m)/(j_1!\cdots j_m!)=(n+1)!/(j_1!\cdots j_m!)$.
This proves the formula by induction.

Apply this exact product derivative identity to the signed parent
equations and then the triangle inequality. For instance
\[
 \partial_{t_c}^{n+1}N_3
  =-dbK_1^2\partial_{t_c}^n A_1
   -dcK_2^3\partial_{t_c}^n\{(h^2+a^2)P_2\},
\]
and its absolute value is bounded by the displayed recurrence for
$N_{n+1}$. This example shows explicitly that neither inherited
diffusion term is omitted. The other seven equations give exactly
\eqref{tgdc:parent_majorants}. Induction in $n$ proves all eight parent
derivative bounds. The identities $R_3=g^2+b^2$,
$a_3=N_3j/K_3$, $b_3=K_3cq$ and $\lambda_3=dK_3^2R_3$ then prove
\eqref{tgdc:coefficient_majorants} for $n\geq1$ by the same product
formula. Their zero-order bounds were already proved in
\eqref{tgdc:proved_boxes}.

Repeatedly differentiate the exact pair to obtain
\eqref{tgdc:exact_derivative_recurrence}. Its componentwise absolute
value gives \eqref{tgdc:amplitude_majorants}, with the zero and first
orders supplied by Theorem~\ref{tgdc:first_derivatives}. Induction
now proves the two amplitude bounds at every order, and the Euclidean
norm bound follows by summing squares. Finally
$\ddot y=(\dot C_3+C_3^2)y$ and direct multiplication of the two by two
matrix give \eqref{tgdc:second_derivative}. Every recurrence uses
finitely many sums at each specified order. Thus each resulting bound
is a determinate expression in the original physical constants and
actual birth data, rather than an unevaluated supremum of a missing
coefficient derivative.
\end{proof}

\subsection{Derivative transport through the specified source clock}

Let $L_s>0$ be the source slot length used for this actual interval,
and define its specified clock by
\[
 \alpha_c=\tau_a/L_s,\qquad
 t_c=t_b+\alpha_c v,\qquad 0\leq v\leq L_s.
\]
The two actual endpoints are retained. Since $\alpha_c$ is constant
for the fixed trajectory and slot, repeated chain differentiation
gives, exactly,
\begin{equation}\label{tgdc:clock}
 \partial_v^ny(t_b+\alpha_cv)
      =\alpha_c^n\partial_{t_c}^ny(t_b+\alpha_cv),\qquad
 |\partial_v^ny|_2\leq\alpha_c^n\sqrt{T_n^2+O_n^2}.
\end{equation}
In a source chart with fixed label $Q$, exponent $h_{\rm src}$,
and $\partial_t v=Q^{-1-h_{\rm src}}$, the derivative contribution
acting on this amplitude alone is consequently
\[
 \partial_t^ny(t_b+\alpha_cv)
  =(\alpha_cQ^{-1-h_{\rm src}})^n
                      \partial_{t_c}^ny(t_b+\alpha_cv).
\]
This identity holds with the stated fixed chart parameters and affine
clock; derivatives of any additional frame, cutoff, or varying
parameter must be included by the ordinary product rule. In particular
neither the clock factor nor the physical frequency has been replaced
by one. These finite derivative estimates establish no uniform
terminal extension for an infinite sequence.
